\fontfam[LM]
\load[emoji]
\typosize[10/12]
\hyperlinks \Blue \Green

\centerline{\typosize[20/24]\bf State-space exploration}
\smallskip
\centerline{\typosize[14/18]\bf Worked examples}
\medskip
\centerline{A~companion to the State-space documentation}
\smallskip
\centerline{\it Matteo Caoduro}
\bigskip
\nonum\notoc
\sec Contents
\maketoc

This volume teaches state-space exploration through worked problems; you
can start here without reading the user guide. Each example starts with a~
problem, then models its states and transitions before applying an operation.
Familiarity with ordinary Clojure functions and collections is assumed.

To try the examples from the project directory, start a~REPL with
\code{clojure -M}; the source directory is on the classpath. Load the library
with:

\begtt \hisyntax{clojure}
(require '[state-space.core :as ss])
\endtt

Evaluate the examples in order in one REPL session, after loading
\code{state-space.core} as shown above. The library supplies traversal or
propagation behavior; ordinary Clojure functions model each problem and
select its answers.

Start with routes through a~building, then explore shortest routes,
alternative routes, and bounded walks. Later groups develop richer states,
weighted searches, context-dependent moves, bottom-up computations, and
finally propagation through a~cycle. Each group starts with a~small use of
its APIs before applying them to more involved problems.

For each problem, pause after reading the question and work out a~state
representation and an approach before continuing to the solution. The
solutions are written as runnable REPL examples.

Many examples use \code{first} after \code{filter} to stop as soon as a~
useful visit is found. Because the search results are lazy, this does not
need to finish exploring the rest of the space. Examples that collect every
answer, such as subset generation, intentionally consume the full result.

\sec Routes, distances, and alternatives

Begin with rooms and passages, then generate moves instead of storing a~graph.
Compare finding one route, measuring distances, and enumerating alternatives.

\secc Which route does depth-first search find first?

\noindent{\bf Problem.} A~building has rooms \code{:lobby},
\code{:stairs}, \code{:landing}, \code{:gallery}, and \code{:exit}. Its
directed passages, in the order inspected from each room, are
\code{:lobby -> :stairs}, \code{:lobby -> :gallery},
\code{:stairs -> :landing}, \code{:landing -> :exit}, and
\code{:gallery -> :exit}. Starting at \code{:lobby}, what route to
\code{:exit} does depth-first search find first?

\smallskip
\noindent{\bf Solution.} Represent each room as a~state and each passage as
a~successor. First inspect the visits returned by DFS:

\begtt \hisyntax{clojure}
(def building
  {:lobby [:stairs :gallery]
   :stairs [:landing]
   :landing [:exit]
   :gallery [:exit]
   :exit []})

(ss/dfs building [:lobby])
;; => ({:state :lobby} {:state :stairs} {:state :landing}
;;     {:state :exit} {:state :gallery})
\endtt

The second argument lists starting states: here there is just one,
\code{:lobby}. Each result is a~visit map containing \code{:state}; the
shared exit is visited only once. The vectors in \code{building} determine
which passage is inspected first. Add \code{path-extension} to record the
route in each visit, then select the first exit visit:

\begtt \hisyntax{clojure}
(first (filter #(= :exit (:state %))
               (ss/dfs building [:lobby]
                           {:extensions [ss/path-extension]})))
;; => {:state :exit, :path [:lobby :stairs :landing :exit]}
\endtt

DFS takes the stairs branch first and returns a~three-passage route, even
though the gallery gives a~shorter route. DFS answers which route it finds
first, not which route is shortest. Change only the search operation to BFS
to find the route with the fewest passages in the same building:

\begtt \hisyntax{clojure}
(first (filter #(= :exit (:state %))
               (ss/bfs building [:lobby]
                           {:extensions [ss/path-extension]})))
;; => {:state :exit, :path [:lobby :gallery :exit]}
\endtt

BFS reaches the exit in two passages. Try reversing the lobby's successors:
DFS now takes the gallery first, while BFS still finds a~shortest route.

\secc Find a~way around a~blocked square

\noindent{\bf Problem.} A~robot is on a $5 \times 5$ grid whose positions
are \code{[x y]} for \code{x} and \code{y} from 0 through 4. It starts at
\code{[0 2]} and can move one square horizontally or vertically. It cannot
leave the grid or enter any of the blocked squares \code{[2 0]},
\code{[2 1]}, \code{[2 2]}, and \code{[2 3]}. What is the fewest moves
needed to reach \code{[4 2]}?

\smallskip
\noindent{\bf Solution.} A~state is the robot's position. Generate adjacent
positions and keep those on the grid and outside the blocked set. BFS
explores states in order of depth, so the depth extension reports the
fewest moves:

\begtt \hisyntax{clojure}
(def blocked-squares #{[2 0] [2 1] [2 2] [2 3]})

(defn grid-moves [[x y]]
  (filter (fn [[nx ny]]
            (and (<= 0 nx 4)
                 (<= 0 ny 4)
                 (not (contains? blocked-squares [nx ny]))))
          [[(inc x) y] [(dec x) y] [x (inc y)] [x (dec y)]]))

(->> (ss/bfs grid-moves [[0 2]]
                 {:extensions [ss/depth-extension]})
     (filter #(= [4 2] (:state %)))
     first
     :depth)
;; => 8
\endtt

This is an implicit graph: the successor function generates legal moves
when a~position is explored instead of storing the whole graph.
\code{[[0 2]]} lists one starting state, the position \code{[0 2]}.
The depth extension starts at zero and counts one per move.

\secc Turn a~word into another word

\noindent{\bf Problem.} The allowed dictionary consists of \code{cold},
\code{cord}, \code{card}, \code{ward}, \code{warm}, \code{worm},
\code{word}, and \code{corm}. Starting from \code{cold}, reach \code{warm}
by changing exactly one letter at a~time. Every word in the route,
including its start and end, must be in the dictionary. Find a~route with
the fewest changes.

\smallskip
\noindent{\bf Solution.} Regard each word as a~state and compute its
neighbours from the dictionary. BFS finds the shortest route:

\begtt \hisyntax{clojure}
(def words ["cold" "cord" "card" "ward" "warm"
            "worm" "word" "corm"])

(defn one-letter-apart? [a b]
  (= 1 (count (filter true? (map not= a b)))))

(defn ladder-neighbours [word]
  (filter #(one-letter-apart? word %) words))

(->> (ss/bfs ladder-neighbours ["cold"]
                 {:extensions [ss/path-extension]})
     (filter #(= "warm" (:state %)))
     first
     :path)
;; => ["cold" "cord" "card" "ward" "warm"]
\endtt

The space is implicit, and lazy \code{filter} stops consuming BFS as soon as
it reaches the goal. The dictionary is an ordered vector, so its order
determines which of several shortest routes BFS returns first.

\secc Find each cell's distance to the nearest zero

\noindent{\bf Problem.} A~grid contains zeros and ones. For every cell,
find its Manhattan distance to the nearest zero. For example, in
\code{["000" "010" "111"]}, the middle cell in the second row is one step
from a~zero, while the bottom-right cell is two steps from one.

\smallskip
\noindent{\bf Solution.} Start one BFS from every zero at once. BFS visits
each cell at its minimum distance from any start; collect those distances
into a~map and arrange them back into rows:

\begtt \hisyntax{clojure}
(def grid ["000" "010" "111"])

(def starts
  (for [row (range (count grid))
        col (range (count (first grid)))
        :when (= \0 (get-in grid [row col]))]
    [row col]))

(defn neighbours [[row col]]
  (filter #(get-in grid %)
          [[(dec row) col] [(inc row) col]
           [row (dec col)] [row (inc col)]]))

(def distance
  (into {} (map (juxt :state :depth))
        (ss/bfs neighbours starts
                    {:extensions [ss/depth-extension]})))

(mapv (fn [row]
        (mapv #(distance [row %]) (range (count (first grid)))))
      (range (count grid)))
;; => [[0 0 0] [0 1 0] [1 2 1]]
\endtt

The multiple starting states make each cell's first BFS arrival its distance
to the nearest zero.

\secc Browse links up to two clicks away

\noindent{\bf Problem.} A~website has these outgoing links, in the order
inspected: from \code{:home} to \code{:about} and \code{:shop}; from
\code{:about} to \code{:team} and \code{:shop}; from \code{:shop} to
\code{:about}, \code{:team}, and \code{:item}; from \code{:team} to
\code{:home} and \code{:reviews}; and from \code{:item} to \code{:reviews}.
The reviews page has no links. Starting at \code{:home}, list every page
reachable by following at most two links, with its minimum number of clicks.
Include \code{:home} itself, but do not follow links from pages reached in
two clicks.

\smallskip
\noindent{\bf Solution.} Use BFS and add the depth extension. Keep the
website as a~map from pages to links; \code{:expand?} prevents requesting
links from visits at depth two, without excluding those pages from the result:

\begtt \hisyntax{clojure}
(def website
  {:home [:about :shop]
   :about [:team :shop]
   :shop [:about :team :item]
   :team [:home :reviews]
   :item [:reviews]
   :reviews []})

(ss/bfs website [:home]
        {:extensions [ss/depth-extension]
         :expand? #(< (:depth %) 2)})
;; => ({:state :home :depth 0} {:state :about :depth 1}
;;     {:state :shop :depth 1} {:state :team :depth 2}
;;     {:state :item :depth 2})
\endtt

Visits are still deduplicated by \code{:state}: BFS guarantees that the first
time a~page is expanded is its shallowest arrival. Pages at the two-click
limit are included, but produce no successors; cycles and repeated links do
not produce duplicate visits.

Alternatively, let the space function read the full visit with
\code{:project identity}. It can inspect \code{:depth} as well as
\code{:state} and return no links at the limit:

\begtt \hisyntax{clojure}
(defn website-links [{:keys [state depth]}]
  (when (< depth 2)
    (website state)))

(ss/bfs website-links [:home]
        {:project identity
         :extensions [ss/depth-extension]})
;; => ({:state :home :depth 0} {:state :about :depth 1}
;;     {:state :shop :depth 1} {:state :team :depth 2}
;;     {:state :item :depth 2})
\endtt

Both versions return the same visits. Prefer \code{:expand?} for this
cutoff; a~search-aware space becomes useful when context changes which
successors are available, as in
\ref[scheduled-connection]{Catch a~scheduled connection} later.

\secc List all routes without repeating a~place

\noindent{\bf Problem.} A~directed map has places \code{:home},
\code{:park}, \code{:market}, and \code{:museum}. Its only roads are
\code{:home -> :park}, \code{:home -> :market},
\code{:park -> :market}, \code{:market -> :park},
\code{:park -> :museum}, and \code{:market -> :museum}. List every route
from \code{:home} to \code{:museum} that visits no place more than once.

\smallskip
\noindent{\bf Solution.} Use \code{simple-paths}: unlike DFS, it keeps
separate routes that reach the same place, while preventing a~route from
looping back to a~place it has already visited. The path extension records
each route prefix:

\begtt \hisyntax{clojure}
(def city-map
  {:home [:park :market]
   :park [:market :museum]
   :market [:park :museum]
   :museum []})

(->> (ss/simple-paths city-map [:home]
                          {:extensions [ss/path-extension]})
     (filter #(= :museum (:state %)))
     (map :path))
;; => ([:home :park :market :museum]
;;     [:home :park :museum]
;;     [:home :market :park :museum]
;;     [:home :market :museum])
\endtt

This enumeration is useful when all routes matter, rather than only the
first route or a~shortest route. Even in a~finite graph, the number of
simple paths can grow quickly.

\secc Find all walks of exactly three steps

\noindent{\bf Problem.} A~directed graph has states \code{:a}, \code{:b},
and \code{:c}, with only the edges \code{:a -> :b}, \code{:b -> :a},
\code{:b -> :c}, and \code{:c -> :a}. List every walk of exactly three
steps starting at \code{:a}. A~walk may visit the same state more than once.

\smallskip
\noindent{\bf Solution.} Stop expanding visits at depth three, but keep
states distinct by route: disable the default global \code{:visit-key}
deduplication and record each path. Then select the visits at the required
depth:

\begtt \hisyntax{clojure}
(def walk-graph
  {:a [:b]
   :b [:a :c]
   :c [:a]})

(->> (ss/bfs walk-graph [:a]
                 {:visit-key nil
                  :expand? #(< (:depth %) 3)
                  :extensions [ss/depth-extension
                               ss/path-extension]})
     (filter #(= 3 (:depth %)))
     (map :path)
     set)
;; => #{[:a :b :a :b] [:a :b :c :a]}
\endtt

This differs from \code{simple-paths}: here revisiting a~state is allowed,
because the question is about walks of a~fixed length, not cycle-free routes.

\secc Check whether a~graph is bipartite

\noindent{\bf Problem.} An undirected graph has two disconnected components:
a~four-cycle and a~triangle. Can its vertices be split into two groups so
that every edge connects different groups?

\smallskip
\noindent{\bf Solution.} In each component, BFS depth parity gives a~two-coloring
candidate. The graph is bipartite exactly when every edge joins vertices of
opposite parity. Start a~new BFS at every still-uncoloured component so a~
disconnected triangle cannot be missed:

\begtt \hisyntax{clojure}
(def graph
  {:a [:b :d] :b [:a :c] :c [:b :d] :d [:a :c]
   :x [:y :z] :y [:x :z] :z [:x :y]})

(def nodes (into (set (keys graph)) (mapcat identity (vals graph))))

(def adjacent
  (reduce (fn [g [from to]]
            (-> g (update from conj to) (update to conj from)))
          (zipmap nodes (repeat []))
          (for [[from tos] graph, to tos] [from to])))

(loop [uncoloured nodes]
  (if-let [root (first uncoloured)]
    (let [component
          (into {} (map (fn [{:keys [state depth]}]
                          [state (mod depth 2)]))
                (ss/bfs adjacent [root]
                            {:extensions [ss/depth-extension]}))]
      (if (some (fn [[node colour]]
                  (some #(= colour (component %)) (adjacent node)))
                component)
        false
        (recur (reduce disj uncoloured (keys component)))))
    true))
;; => false
\endtt

The square is bipartite, but the separate triangle contains an edge whose
endpoints have the same parity. Checking every component is essential.

\sec Modeling choices and richer states

Keep the familiar DFS and BFS calls while changing what a~state represents.
The problems move from partial choices and puzzle configurations to progress,
resources, and interacting processes.

\secc Generate every subset

\noindent{\bf Problem.} List every subset of the four distinct items
\code{:a}, \code{:b}, \code{:c}, and \code{:d}, including the empty set and
the set containing all four. Represent each subset as a~vector in the
items' original order.

\smallskip
\noindent{\bf Solution.} At each step, decide whether to leave the next
item out or add it to the partial result. DFS explores this decision tree;
the index in the state tells us when all decisions have been made:

\begtt \hisyntax{clojure}
(def items [:a :b :c :d])

(defn subset-decisions [[i chosen]]
  (when (< i (count items))
    [[(inc i) chosen]
     [(inc i) (conj chosen (items i))]]))

(->> (ss/dfs subset-decisions [[0 []]])
     (map :state)
     (filter #(= (count items) (first %)))
     (map second))
;; => ([] [:d] [:c] [:c :d]
;;     [:b] [:b :d] [:b :c] [:b :c :d]
;;     [:a] [:a :d] [:a :c] [:a :c :d]
;;     [:a :b] [:a :b :d] [:a :b :c] [:a :b :c :d])
\endtt

The search state includes both the partial answer and how many choices have
been considered. This lets DFS distinguish unfinished decisions from
completed subsets.

\secc Place eight queens without conflicts

\noindent{\bf Problem.} Place eight queens on a~chessboard so that no two
share a~row, column, or diagonal. How many solutions are there?

\smallskip
\noindent{\bf Solution.} Place one queen per row. A~partial placement is a~
vector of column numbers, so row conflicts are impossible by construction.
Reject a~column as soon as it conflicts with a~queen already placed; DFS
then explores only safe partial boards:

\begtt \hisyntax{clojure}
(defn safe-column? [placed column]
  (let [row (count placed)]
    (every? (fn [[previous-row previous-column]]
              (and (not= column previous-column)
                   (not= (- row previous-row)
                         (Math/abs (- column previous-column)))))
            (map-indexed vector placed))))

(defn queen-placements [placed]
  (when (< (count placed) 8)
    (for [column (range 8)
          :when (safe-column? placed column)]
      (conj placed column))))

(count (filter #(= 8 (count (:state %)))
               (ss/dfs queen-placements [[]])))
;; => 92
\endtt

The partial column vector is the search state. Pruning unsafe choices as
soon as they appear avoids exploring complete boards that could never be
solutions.

\secc Make a~number from the fewest perfect squares

\noindent{\bf Problem.} Starting from 12, repeatedly subtract a~positive
perfect square no larger than the current number. What is the fewest
subtractions needed to reach zero?

\smallskip
\noindent{\bf Solution.} Use the remaining number as the state, with one
successor for each legal square subtraction. BFS finds a~route with the
fewest subtractions:

\begtt \hisyntax{clojure}
(defn square-subtractions [remaining]
  (for [root (range 1 (inc (long (Math/sqrt remaining))))]
    (- remaining (* root root))))

(->> (ss/bfs square-subtractions [12]
                 {:extensions [ss/path-extension]})
     (filter #(zero? (:state %)))
     first
     :path)
;; => [12 8 4 0]
\endtt

Each step subtracts a~square; the path therefore represents
\code{12 = 4 + 4 + 4}.

\secc Fewest moves across a~chessboard

\noindent{\bf Problem.} A~knight starts at \code{[0 0]} on a~standard
$8 \times 8$ chessboard whose squares are \code{[x y]} with coordinates
from 0 through 7. A~knight moves two squares in one direction and one square
perpendicular to that direction, staying on the board. What is the fewest
moves needed to reach the opposite corner, \code{[7 7]}?

\smallskip
\noindent{\bf Solution.} Treat each square as a~state. Generate the legal
knight moves from a~square, then use BFS because every move has equal cost.

\begtt \hisyntax{clojure}
(defn knight-moves [[x y]]
  (for [[dx dy] [[1 2] [2 1] [-1 2] [-2 1]
                 [1 -2] [2 -1] [-1 -2] [-2 -1]]
        :let [next [(+ x dx) (+ y dy)]]
        :when (every? #(<= 0 % 7) next)]
    next))

(->> (ss/bfs knight-moves [[0 0]]
                 {:extensions [ss/depth-extension]})
     (filter #(= [7 7] (:state %)))
     first
     :depth)
;; => 6
\endtt

The successor function generates only legal squares; no graph needs to be
stored.

\secc[jug-puzzle] Measure four litres with two jugs

\noindent{\bf Problem.} You have an unmarked three-litre jug and an unmarked
five-litre jug, both initially empty. You may fill either jug completely
from a~tap, empty either jug, or pour from one jug to the other until the
source is empty or the destination is full. What is the fewest actions
needed to leave exactly four litres in the five-litre jug?

\smallskip
\noindent{\bf Solution.} Represent a~state as the amount of water in both
jugs. Each allowed action creates a~successor state. BFS finds a~solution
with the fewest actions, and the path extension shows the states along it:

\begtt \hisyntax{clojure}
(defn jug-moves [[a b]]
  (let [capacity-a 3 capacity-b 5
        pour-ab (min a (- capacity-b b))
        pour-ba (min b (- capacity-a a))]
    [[capacity-a b] [a capacity-b] [0 b] [a 0]
     [(- a pour-ab) (+ b pour-ab)]
     [(+ a pour-ba) (- b pour-ba)]]))

(->> (ss/bfs jug-moves [[0 0]]
                 {:extensions [ss/path-extension]})
     (filter (fn [{[_ b] :state}] (= 4 b)))
     first
     :path)
;; => [[0 0] [0 5] [3 2] [0 2] [2 0] [2 5] [3 4]]
\endtt

The full state matters: the same amount in one jug can lead to different
future moves depending on the other jug's contents.

\secc Help the missionaries and cannibals cross

\noindent{\bf Problem.} Three missionaries and three cannibals are on one
bank of a~river. A~boat carries one or two people. Whenever missionaries
are present on a~bank, they must not be outnumbered by cannibals. What is
the fewest crossings that take everyone to the opposite bank?

\smallskip
\noindent{\bf Solution.} Record how many missionaries and cannibals remain
on the starting bank, plus which bank has the boat. A~crossing is legal only
when both banks remain safe. BFS finds the minimum number of crossings:

\begtt \hisyntax{clojure}
(defn safe-bank? [missionaries cannibals]
  (or (zero? missionaries)
      (>= missionaries cannibals)))

(defn safe-state? [[left-m left-c _boat]]
  (and (<= 0 left-m 3)
       (<= 0 left-c 3)
       (safe-bank? left-m left-c)
       (safe-bank? (- 3 left-m) (- 3 left-c))))

(defn crossings [[missionaries cannibals boat]]
  (let [direction (if (zero? boat) -1 1)]
    (for [[move-m move-c] [[1 0] [2 0] [0 1] [0 2] [1 1]]
          :let [next [(+ missionaries (* direction move-m))
                      (+ cannibals (* direction move-c))
                      (- 1 boat)]]
          :when (safe-state? next)]
      next)))

(->> (ss/bfs crossings [[3 3 0]]
                 {:extensions [ss/depth-extension]})
     (filter #(= [0 0 1] (:state %)))
     first
     :depth)
;; => 11
\endtt

Checking both banks after every move keeps unsafe states out of the search.
The boat side is part of the state because it determines the direction of
the next crossing.

\secc Open a~combination lock in the fewest turns

\noindent{\bf Problem.} A~combination lock has four decimal wheels,
initially set to \code{"0000"}. One turn moves one wheel up or down by one,
wrapping between 0 and 9. The combinations \code{"0201"}, \code{"0101"},
\code{"0102"}, \code{"1212"}, and \code{"2002"} are blocked and cannot be
entered. What is the fewest turns needed to reach \code{"0202"}?

\smallskip
\noindent{\bf Solution.} A~four-character string is a~state. Generate the
eight possible single-wheel turns, wrapping each digit with modular
arithmetic, and omit blocked combinations. BFS finds the minimum number of
turns:

\begtt \hisyntax{clojure}
(def blocked-codes #{"0201" "0101" "0102" "1212" "2002"})

(defn lock-turns [code]
  (for [wheel (range 4)
        direction [-1 1]
        :let [digit (Character/digit ^char (nth code wheel) 10)
              next-digit (mod (+ digit direction) 10)
              next-code (str (subs code 0 wheel)
                             next-digit
                             (subs code (inc wheel)))]
        :when (not (blocked-codes next-code))]
    next-code))

(->> (ss/bfs lock-turns ["0000"]
                 {:extensions [ss/depth-extension]})
     (filter #(= "0202" (:state %)))
     first
     :depth)
;; => 6
\endtt

The search need not construct all 10,000 combinations in advance; it
generates legal turns only for combinations it explores.

\secc Solve a~sliding-tile puzzle

\noindent{\bf Problem.} A~two-row, three-column board contains tiles
\code{1} through \code{5} and one blank, represented by \code{0}. Represent
the board as a~vector in row-major order. The blank may swap places with a~
tile directly above, below, left, or right of it, provided that position is
on the board. Starting from \code{[4 1 3 0 2 5]}, find the fewest swaps
that produce \code{[1 2 3 4 5 0]}.

\smallskip
\noindent{\bf Solution.} Represent the board as a~vector in row-major
order. A~state is the whole board; generate successors by swapping the
blank with each adjacent position. BFS finds a~shortest solution:

\begtt \hisyntax{clojure}
(defn slide-blank [tiles]
  (let [blank (.indexOf ^java.util.List tiles 0)
        row (quot blank 3)
        col (rem blank 3)]
    (for [[dr dc] [[-1 0] [1 0] [0 -1] [0 1]]
          :let [r (+ row dr)
                c (+ col dc)]
          :when (and (<= 0 r 1)
                     (<= 0 c 2))]
      (assoc tiles blank
             (tiles (+ (* r 3) c))
             (+ (* r 3) c) 0))))

(->> (ss/bfs slide-blank [[4 1 3 0 2 5]]
                 {:extensions [ss/path-extension]})
     (filter #(= [1 2 3 4 5 0] (:state %)))
     first
     :path)
;; => [[4 1 3 0 2 5]
;;     [0 1 3 4 2 5]
;;     [1 0 3 4 2 5]
;;     [1 2 3 4 0 5]
;;     [1 2 3 4 5 0]]
\endtt

The path contains four swaps. Unlike the \ref[jug-puzzle]{jug puzzle}, here
each move changes the arrangement of the whole state rather than quantities
in separate containers.

\secc Push a~box to its goal

\noindent{\bf Problem.} In the small warehouse below, move the box marked
\code{\$} onto the goal marked \code{.}. The player may walk through empty
cells and push the box, but cannot move through walls or pull the box. What
is the fewest player moves?

\begtt
#######
#@    #
# $ . #
#     #
#######
\endtt

\smallskip
\noindent{\bf Solution.} A~state contains both the player's position and
the box's position. When the player steps into the box, that same move
pushes it one cell farther; reject the move if the box would hit a~wall.
BFS searches these joint configurations by number of player moves:

\begtt \hisyntax{clojure}
(def warehouse
  ["#######"
   "#@    #"
   "# $ . #"
   "#     #"
   "#######"])

(defn wall? [[row column]]
  (= \# (get-in warehouse [row column])))

(defn sokoban-moves [[player box]]
  (for [direction [[1 0] [-1 0] [0 1] [0 -1]]
        :let [next-player (mapv + player direction)]
        :when (not (wall? next-player))
        :let [next-box (if (= next-player box)
                         (mapv + box direction)
                         box)]
        :when (not (wall? next-box))]
    [next-player next-box]))

(->> (ss/bfs sokoban-moves [[[1 1] [2 2]]]
                 {:extensions [ss/depth-extension]})
     (filter #(= [2 4] (second (:state %))))
     first
     :depth)
;; => 3
\endtt

Combining the two positions in one state lets the ordinary visited-state
check distinguish configurations with the box in different places.

\secc Plan around the Sussman anomaly

\noindent{\bf Problem.} Three blocks are arranged in two stacks, with the
bottom block first: \code{[:a :c]} and \code{[:b]}. Move one clear block at
a~time, either onto the table or onto another clear block, until the stack
\code{[:c :b :a]} is built. What is the fewest moves?

\smallskip
\noindent{\bf Solution.} Represent a~world as a~set of stacks. A~move
removes a~stack's top block and places it on the table or on another stack's
top. BFS searches complete world states, so it can find a~plan that first
undoes part of the initial arrangement:

\begtt \hisyntax{clojure}
(defn block-moves [stacks]
  (for [from stacks
        :let [block (peek from)
              without (cond-> (disj stacks from)
                        (seq (pop from)) (conj (pop from)))]
        target (cons nil (disj stacks from))]
    (if target
      (-> without (disj target) (conj (conj target block)))
      (conj without [block]))))

(let [plan (first (filter #(= #{[:c :b :a]} (:state %))
                          (ss/bfs block-moves [#{[:a :c] [:b]}]
                                      {:extensions [ss/depth-extension
                                                    ss/path-extension]})))]
  {:depth (:depth plan)
   :path (mapv #(vec (sort-by first %)) (:path plan))})
;; => {:depth 3, :path [[[:a :c] [:b]] [[:a] [:b] [:c]]
;;                      [[:a] [:c :b]] [[:c :b :a]]]}
\endtt

The planner must temporarily move \code{:c} off \code{:a} so that it can
build \code{:b} on \code{:c}, then put \code{:a} on top. Treating the goals
“put \code{:c} on the table” and “put \code{:a} on \code{:b}” as independent
subtasks would miss that interaction. The initial world is a~set, so its
iteration order can select a~different shortest plan on another runtime.

\secc Visit every node in the fewest steps

\noindent{\bf Problem.} A~robot moves along an undirected star-shaped graph:
node 0 is connected to nodes 1, 2, and 3, with no other edges. It may start
at any node and revisit nodes. What is the fewest moves needed to visit all
four nodes?

\smallskip
\noindent{\bf Solution.} The current node alone is not enough to describe
the search: future success also depends on which nodes have already been
visited. Use the pair \code{[node visited]} as the state and start a~BFS
from each possible node:

\begtt \hisyntax{clojure}
(def graph [[1 2 3]
            [0]
            [0]
            [0]])

(def all-nodes (set (range (count graph))))

(defn visit-next [[node visited]]
  (for [next (graph node)]
    [next (conj visited next)]))

(def starts
  (map (fn [node] [node #{node}])
       (range (count graph))))

(->> (ss/bfs visit-next starts
                 {:extensions [ss/depth-extension]})
     (filter #(= all-nodes (second (:state %))))
     first
     :depth)
;; => 4
\endtt

Because the visited set is part of state identity, reaching the same graph
node with different progress is not mistaken for a~duplicate visit.

\secc Collect samples and return to base

\noindent{\bf Problem.} A~small rover moves horizontally or vertically on a~
three-by-three grid with coordinates from 0 through 2. It starts at
\code{[0 0]} with four units of fuel; each move uses one unit. It must
collect samples at \code{[2 0]} and \code{[2 2]} by visiting those cells,
then return to its starting position. At \code{[1 1]} it may refill to four
units. Each grid move and each refuelling action takes one step; collecting
a~sample takes no extra step. What is the fewest steps?

\smallskip
\noindent{\bf Solution.} Position alone is not a~complete state: fuel and
collected samples determine which actions remain possible and whether the
goal has been reached. Store all three in each state. The default visit
identity then distinguishes configurations that share a~position but differ
in fuel or samples:

\begtt \hisyntax{clojure}
(def samples {[2 0] :amber, [2 2] :cobalt})
(def station [1 1])
(def base [0 0])
(def capacity 4)

(defn rover-moves [{:keys [position fuel collected]}]
  (let [moves
        (for [[dx dy] [[-1 0] [1 0] [0 -1] [0 1]]
              :let [next [(+ (first position) dx)
                          (+ (second position) dy)]]
              :when (and (every? #(<= 0 % 2) next) (pos? fuel))]
          {:position next
           :fuel (dec fuel)
           :collected (cond-> collected
                        (samples next) (conj (samples next)))})]
    (if (and (= position station) (< fuel capacity))
      (conj (vec moves)
            {:position position
             :fuel capacity
             :collected collected})
      moves)))

(def initial-rover
  {:position base :fuel capacity :collected #{}})

(->> (ss/bfs rover-moves [initial-rover]
                 {:extensions [ss/depth-extension]})
     (filter #(and (= base (get-in % [:state :position]))
                   (= #{:amber :cobalt}
                      (get-in % [:state :collected]))))
     first
     :depth)
;; => 12
\endtt

BFS finds the minimum number of actions because moves and refuelling each
count as one step. This is a~product state space: a~physical position is
combined with resource and collection progress.

\secc Check mutual exclusion across process interleavings

\noindent{\bf Problem.} Two processes use Peterson's mutual-exclusion
protocol. Can both enter their critical sections at once? What shortest
counterexample appears if the waiting condition is removed?

\smallskip
\noindent{\bf Solution.} A~state records each process's control location,
its request flag, and whose turn it is. At each step either process may
advance by one instruction. BFS can check that the correct protocol never
reaches a~state where both program counters are \code{3}, then return a~
shortest trace when the waiting condition is broken:

\begtt \hisyntax{clojure}
(def initial-system {:pc [0 0] :flag [false false] :turn 0})

(defn process-steps [wait?]
  (fn [{:keys [turn] :as system}]
    (for [i [0 1]
          :let [j (- 1 i)]]
      (case (get-in system [:pc i])
        0 (-> system (assoc-in [:flag i] true) (assoc-in [:pc i] 1))
        1 (-> system (assoc :turn j) (assoc-in [:pc i] 2))
        2 (if (and wait? (get-in system [:flag j]) (= turn j))
            system
            (assoc-in system [:pc i] 3))
        3 (assoc-in system [:pc i] 4)
        4 (-> system (assoc-in [:flag i] false)
              (assoc-in [:pc i] 0))))))

(defn violation [wait?]
  (first (filter #(= [3 3] (get-in % [:state :pc]))
                 (ss/bfs (process-steps wait?) [initial-system]
                             {:extensions [ss/depth-extension
                                           ss/path-extension]}))))

(let [counterexample (violation false)]
  [(nil? (violation true))
   (:depth counterexample)
   (mapv #(get-in % [:pc]) (:path counterexample))])
;; => [true 6 [[0 0] [1 0] [2 0] [3 0] [3 1] [3 2] [3 3]]]
\endtt

The correct protocol preserves mutual exclusion for every interleaving.
Without waiting, the first process enters and stays in its critical
section while the second also enters, violating the safety property after
six instructions.

\sec Weighted routes and custom priorities

Moves no longer all cost the same. Start with travel time and edge details,
then change the objective or guide exploration with an estimate.

\secc Choose the least expensive route

\noindent{\bf Problem.} A~travel network has directed trips with these
times in hours: from \code{:home} to \code{:museum} (11), \code{:village}
(4), or \code{:park} (1); from \code{:village} to \code{:museum} (7) or
\code{:hill} (2); from \code{:park} to \code{:museum} (8) or
\code{:village} (1); and from \code{:hill} to \code{:museum} (2). There
are no other trips. Which route from \code{:home} to \code{:museum} takes
the least time, and what is its total cost?

\smallskip
\noindent{\bf Solution.} Each edge stores its travel time. Decode an edge
into a~transition containing both its destination and weight; best-first
uses the accumulated distance to return the least-cost route:

\begtt \hisyntax{clojure}
(def travel-times
  {:home [[:museum 11] [:village 4] [:park 1]]
   :village [[:museum 7] [:hill 2]]
   :park [[:museum 8] [:village 1]]
   :hill [[:museum 2]]
   :museum []})

(->> (ss/best-first travel-times [:home]
                        {:decode-edge (fn [[state weight]]
                                        {:state state :weight weight})
                         :extensions [ss/path-extension]})
     (filter #(= :museum (:state %)))
     first)
;; => {:state :museum
;;     :path [:home :park :village :hill :museum], :distance 6}
\endtt

Unlike BFS, best-first minimizes the sum of edge costs, not the number of
edges. Its default distance priority is Dijkstra's algorithm when edge
weights are non-negative.

\secc Keep the details of each leg

\noindent{\bf Problem.} A~travel network has these directed trips:
\code{:home -> :museum} by express train (11~hours), \code{:home ->
:village} by ferry (4), \code{:home -> :park} by tram (1),
\code{:village -> :museum} by bus (7), \code{:village -> :hill} by coach
(2), \code{:park -> :museum} by taxi (8), \code{:park -> :village} on
foot (1), and \code{:hill -> :museum} by rail (2). There are no other
trips. Find the least-time route from \code{:home} to \code{:museum}, and
report each leg's mode and time as well as the total time.

\smallskip
\noindent{\bf Solution.} Give each edge its destination, time, and mode.
Decode these into transition maps and use \code{trail-extension} to retain
the full transitions along the selected route:

\begtt \hisyntax{clojure}
(def labelled-travel
  {:home [[:museum 11 :express]
          [:village 4 :ferry]
          [:park 1 :tram]]
   :village [[:museum 7 :bus]
             [:hill 2 :coach]]
   :park [[:museum 8 :taxi]
          [:village 1 :foot]]
   :hill [[:museum 2 :rail]]
   :museum []})

(->> (ss/best-first labelled-travel [:home]
                        {:decode-edge (fn [[state weight mode]]
                                        {:state state
                                         :weight weight
                                         :mode mode})
                         :extensions [ss/trail-extension]})
     (filter #(= :museum (:state %)))
     first)
;; => {:state :museum
;;     :trail [{:state :home}
;;             {:state :park, :weight 1, :mode :tram}
;;             {:state :village, :weight 1, :mode :foot}
;;             {:state :hill, :weight 2, :mode :coach}
;;             {:state :museum, :weight 2, :mode :rail}],
;;     :distance 6}
\endtt

The trail preserves edge data that a~state-only path would lose.

\secc Two travelers meet as quickly as possible

\noindent{\bf Problem.} Two travelers start in \code{:west} and
\code{:east}. On each turn, both must move along one road. A~turn lasts as
long as the slower traveler takes to reach the next city; the faster traveler
waits. Find the minimum elapsed time until they are in the same city.

\smallskip
\noindent{\bf Solution.} A~state is the pair of current cities. Each
successor is one pair of simultaneous moves, and its cost is the longer of
the two road times. Dijkstra search minimizes the sum of these per-turn
durations. This models the joint state directly rather than searching for
each traveler independently:

\begtt \hisyntax{clojure}
(def meeting-roads
  {:west [[:ridge 2]]
   :ridge [[:west 2] [:town 2] [:bridge 5]]
   :east [[:market 1]]
   :market [[:east 1] [:town 2] [:bridge 4]]
   :town [[:ridge 2] [:market 2]]
   :bridge [[:ridge 5] [:market 4]]})

(defn simultaneous-moves [[city-a city-b]]
  (for [[next-a time-a] (meeting-roads city-a)
        [next-b time-b] (meeting-roads city-b)]
    {:state [next-a next-b]
     :weight (max time-a time-b)}))

(->> (ss/best-first simultaneous-moves
                        [[:west :east]]
                        {:decode-edge identity
                         :extensions [ss/path-extension]})
     (filter (fn [{[city-a city-b] :state}]
               (= city-a city-b)))
     first)
;; => {:state [:town :town]
;;     :path [[:west :east] [:ridge :market] [:town :town]]
;;     :distance 4}
\endtt

The two road times are not added within a~turn: one traveler waits for the
other. Modeling the pair as one state and the elapsed turn duration as one
edge cost makes the ordinary weighted search objective match the problem.

\secc Make a~short presentation covering required topics

\noindent{\bf Problem.} A~short presentation must cover \code{:welcome},
\code{:robots}, and \code{:navigation}. Available clips have a~duration and
cover one or more topics. Find a~presentation with minimum total duration.
This is a~small set-cover search.

\smallskip
\noindent{\bf Solution.} A~state records the topics still to cover and the
clips already selected. At each step choose the first remaining topic and
consider clips that cover it. Each transition's weight is the clip duration,
so best-first returns the least-duration complete presentation:

\begtt \hisyntax{clojure}
(def required-topics #{:welcome :robots :navigation})
(def clips
  [{:id :short-welcome :duration 2 :topics #{:welcome}}
   {:id :robot-navigation :duration 5 :topics #{:robots :navigation}}
   {:id :robot-basics :duration 4 :topics #{:robots}}
   {:id :navigation-demo :duration 2 :topics #{:navigation}}
   {:id :full-introduction :duration 8 :topics #{:welcome :robots}}])

(defn add-clip [{:keys [remaining selected]}]
  (when-let [topic (first (sort remaining))]
    (for [{:keys [id duration topics]} clips
          :when (contains? topics topic)]
      {:state {:remaining (reduce disj remaining topics)
               :selected (conj selected id)}
       :weight duration})))

(def presentation
  (first (filter #(empty? (get-in % [:state :remaining]))
                 (ss/best-first
                  add-clip
                  [{:remaining required-topics :selected []}]
                  {:decode-edge identity}))))

presentation
;; => {:state {:remaining #{}
;;             :selected [:robot-navigation :short-welcome]}
;;     :distance 7}
\endtt

The selected topic makes the branching finite and ensures that every chosen
clip covers at least one outstanding requirement. Best-first compares total
duration, rather than clip count; here one multi-topic clip makes the
seven-minute solution shorter than choosing separate clips for each topic.

\secc Grow a~minimum spanning tree

\noindent{\bf Problem.} Connect four towns using the least total road length.
The undirected roads have lengths \code{:a-:b} (1), \code{:a-:c} (4),
\code{:b-:c} (2), \code{:b-:d} (5), and \code{:c-:d} (1). Which roads
belong to a~minimum spanning tree, and what is their total length?

\smallskip
\noindent{\bf Solution.} Prim's algorithm repeatedly chooses the lightest
edge leaving the already reached towns. Give each visit the weight of the
edge that joins it, and use that value as best-first's priority. This is
not Dijkstra's algorithm: the priority is the joining edge, not the
cumulative distance from the start. The custom extension's \code{:start}
initializes the root's priority; \code{:next} reads the parent and decoded
transition to record the selected road. Both return fields to add to the
visit; the unused visit arguments are written as \code{_}:

\begtt \hisyntax{clojure}
(def roads
  {:a [{:state :b :weight 1} {:state :c :weight 4}]
   :b [{:state :a :weight 1} {:state :c :weight 2}
       {:state :d :weight 5}]
   :c [{:state :a :weight 4} {:state :b :weight 2}
       {:state :d :weight 1}]
   :d [{:state :b :weight 5} {:state :c :weight 1}]})

(def joining-edge
  {:start (fn [_] {:joining-edge 0})
   :next (fn [parent transition _]
           {:joining-edge (:weight transition)
            :road [(:state parent) (:state transition)
                   (:weight transition)]})})

(let [tree (ss/best-first roads [:a]
                              {:decode-edge identity
                               :extensions [joining-edge]
                               :priority :joining-edge})]
  [(mapv :road (rest tree))
   (transduce (map :joining-edge) + tree)])
;; => [[[:a :b 1] [:b :c 2] [:c :d 1]] 4]
\endtt

Each newly reached town is attached through its cheapest available edge.
The three selected roads have total length four.

\secc Find a~route minimizing its largest climb

\noindent{\bf Problem.} A~hiker crosses this height grid from the upper-left
corner to the lower-right, moving one cell horizontally or vertically at
each step. The effort of a~route is its largest height difference between
consecutive cells. Find a~route with the smallest possible effort:

\begtt \hisyntax{clojure}
(def heights [[1 2 2]
              [3 8 2]
              [5 3 5]])
\endtt

\smallskip
\noindent{\bf Solution.} Unlike ordinary route length, the effort is not
the sum of the climbs. Accumulate the largest climb seen so far and let
best-first prioritize that value. \code{ss/accumulate} builds the extension:
start \code{:effort} at zero, then update it from the parent's effort and
the next transition's climb:

\begtt \hisyntax{clojure}
(defn grid-climbs [[row col]]
  (for [[next-row next-col]
        [[(dec row) col] [(inc row) col]
         [row (dec col)] [row (inc col)]]
        :let [here (get-in heights [row col])
              there (get-in heights [next-row next-col])]
        :when there]
    {:state [next-row next-col]
     :climb (abs (- there here))}))

(def effort-extension
  (ss/accumulate :effort (constantly 0)
                     (fn [effort transition]
                       (max effort (:climb transition)))))

(->> (ss/best-first grid-climbs [[0 0]]
                        {:decode-edge identity
                         :extensions [effort-extension]
                         :priority :effort})
     (filter #(= [2 2] (:state %)))
     first)
;; => {:state [2 2], :effort 2}
\endtt

Along any route, adding a~climb can never lower its maximum, so this
priority orders visits by the best effort found so far.

\secc Guide best-first search with an estimate

\noindent{\bf Problem.} From \code{:home}, a~delivery can reach
\code{:hill} in 1~hour, \code{:river} in 2, or \code{:market} in 4.
Each of those places has one route to \code{:goal}, taking 5, 2, and 2~hours
respectively. There are no other routes. The estimated remaining times are
3~hours from home, 4 from the hill, 2 from the river, 2 from the market, and
0 from the goal. Find the least-time route from home to the goal.

\smallskip
\noindent{\bf Solution.} Best-first can prioritize each visit by the time
already travelled plus the estimated time remaining. This is A* search; the
estimates here are consistent, so the first settled goal has minimum travel
time:

\begtt \hisyntax{clojure}
(def delivery-routes
  {:home [[:hill 1] [:river 2] [:market 4]]
   :hill [[:goal 5]]
   :river [[:goal 2]]
   :market [[:goal 2]]
   :goal []})
(def remaining-time
  {:home 3 :hill 4 :river 2 :market 2 :goal 0})

(->> (ss/best-first delivery-routes [:home]
                        {:decode-edge (fn [[state weight]]
                                        {:state state :weight weight})
                         :extensions [ss/path-extension
                                      ss/distance-extension]
                         :priority (fn [{:keys [state distance]}]
                                     (+ distance
                                        (remaining-time state)))})
     (filter #(= :goal (:state %)))
     first)
;; => {:state :goal, :path [:home :river :goal], :distance 4}
\endtt

The estimate guides the order of exploration; it does not change the route
cost. For this once-per-state best-first search, use a~consistent estimate
that is zero at the goal: along each edge, the estimate at its source must
not exceed the edge's time plus the destination's estimate. For example,
home's estimate of 3 is no greater than the 2 hours to the river plus its
estimate of 2. The explicit \code{distance-extension} keeps calculating
travel time when a~custom \code{:priority} replaces the default distance
priority.

\sec Routes with context and constraints

Now the moves available depend on arrival time, fuel, or progress through
a~label pattern. Compare carrying that information in visits with making it
part of the state itself.

\secc[scheduled-connection] Catch a~scheduled connection

\noindent{\bf Problem.} A~traveler leaves \code{:home} at 8:00. A~bus
reaches \code{:junction} at 8:10 and another reaches \code{:museum} directly
at 8:25. From the junction, a~bus leaves at 8:15 and reaches the museum at
8:20. Find the earliest arrival. A~connection can be used only if the
traveler reaches its departure stop in time.

\smallskip
\noindent{\bf Solution.} Store departure and arrival times relative to 8:00.
The space receives the full visit and offers only trips that have not
departed yet. Its transition weight is the wait plus travel time, so the
accumulated Dijkstra distance is the arrival time:

\begtt \hisyntax{clojure}
(def timetable
  {:home [{:to :junction :departs 0 :arrives 10}
          {:to :museum :departs 0 :arrives 25}]
   :junction [{:to :museum :departs 15 :arrives 20}]
   :museum []})

(defn catchable-trips [{:keys [state distance]}]
  (for [{:keys [to departs arrives]} (timetable state)
        :when (<= distance departs)]
    [to (- arrives distance)]))

(->> (ss/best-first catchable-trips [:home]
                        {:project identity
                         :decode-edge (fn [[state weight]]
                                        {:state state :weight weight})
                         :extensions [ss/path-extension]})
     (filter #(= :museum (:state %)))
     first)
;; => {:state :museum, :path [:home :junction :museum], :distance 20}
\endtt

The scheduled edges are search-aware: whether a~trip is a~successor depends
on the current arrival time in \code{:distance}. The default
\code{:visit-key} remains \code{:state}; for a~timetable where earlier arrival
always leaves at least as many future options, a~later visit to the same stop
is dominated. A~static stop-to-stop graph cannot directly express which
departures are still catchable; without search-aware generation, the time
would need to be encoded into states or represented with a~time-expanded
graph.

\secc Reach the goal with limited fuel

\noindent{\bf Problem.} A~vehicle starts at \code{:home} with 5~units of
fuel and cannot refuel. Its only directed routes, annotated with fuel used,
are
\code{:home -> :junction} (4), \code{:home -> :detour} (1),
\code{:detour -> :junction} (1), and \code{:junction -> :goal} (3).
Can it reach \code{:goal} without any move using more fuel than remains?
Arriving with zero fuel is allowed.

\smallskip
\noindent{\bf Solution.} Fuel remaining affects which routes are possible,
so the space function needs the full visit, not only its location. Track
fuel with an extension and include it in \code{:visit-key}: arriving at the
same junction with different fuel is a~different search position.

\begtt \hisyntax{clojure}
(def fuel-routes
  {:home [[:junction 4] [:detour 1]]
   :detour [[:junction 1]]
   :junction [[:goal 3]]
   :goal []})

(def fuel-left
  (ss/accumulate :fuel-left (constantly 5)
                     (fn [fuel transition]
                       (- fuel (:fuel-used transition)))))

(defn affordable-routes [{:keys [state fuel-left]}]
  (filter (fn [[_ fuel-used]] (<= fuel-used fuel-left))
          (fuel-routes state)))

(->> (ss/bfs affordable-routes [:home]
                 {:project identity
                  :visit-key (juxt :state :fuel-left)
                  :decode-edge (fn [[state fuel-used]]
                                 {:state state :fuel-used fuel-used})
                  :extensions [fuel-left ss/path-extension]})
     (filter #(= :goal (:state %)))
     first)
;; => {:state :goal, :path [:home :detour :junction :goal], :fuel-left 0}
\endtt

The direct arrival at \code{:junction} leaves too little fuel to continue.
If visits were identified only by location, BFS would discard the later
arrival with enough fuel. Including \code{:fuel-left} in \code{:visit-key}
preserves both possibilities.

\secc Find a~route matching a~label pattern

\noindent{\bf Problem.} A~directed network has edges
\code{:start -> :middle} labeled \code{a},
\code{:start -> :wrong} labeled \code{b},
\code{:middle -> :target} labeled \code{b},
\code{:target -> :middle} labeled \code{a}, and
\code{:wrong -> :target} labeled \code{a}, with no other edges. Find the
shortest route from \code{:start} to \code{:target} whose edge labels match
one or more repetitions of the pattern \code{ab}.

\smallskip
\noindent{\bf Solution.} Carry both the current graph node and the state
of a~small finite automaton in the search state. The automaton accepts after
a~\code{b} follows an \code{a}; from that accepting state, another \code{a}
can begin the next repetition. Thus ordinary BFS explores the product of the
graph and the pattern recognizer, while the default visit identity keeps
distinct automaton states at the same graph node separate:

\begtt \hisyntax{clojure}
(def labelled-graph
  {:start [{:to :middle :label "a"}
           {:to :wrong :label "b"}]
   :middle [{:to :target :label "b"}]
   :target [{:to :middle :label "a"}]
   :wrong [{:to :target :label "a"}]})

(defn matching-edges [[node phase]]
  (for [{:keys [to label]} (labelled-graph node)
        :let [next-phase
              (case [phase label]
                [:need-a "a"] :need-b
                [:need-b "b"] :accept
                [:accept "a"] :need-b
                nil)]
        :when next-phase]
    {:state [to next-phase] :label label}))

(let [route (->> (ss/bfs matching-edges [[:start :need-a]]
                             {:decode-edge identity
                              :extensions [ss/trail-extension]})
                (filter #(= [:target :accept] (:state %)))
                first)]
  (mapv :label (rest (:trail route))))
;; => ["a" "b"]
\endtt

This expresses a~regular constraint on one route. It is not a~declarative
graph-pattern query: matching patterns that bind several nodes and
relationships together or impose subgraph-matching semantics may require a~
different abstraction.

\sec Computing after successors

Switch from finding routes to using completed successor results. Begin
with dependency order, then compute bottom-up values before controlling
successor processing and inspecting cycles.

\secc Build tasks in dependency order

\noindent{\bf Problem.} An application has four tasks: \code{:app} needs
\code{:lib} and \code{:config}; each of those needs \code{:base}; and
\code{:base} has no prerequisites. Find an order in which to perform all
four tasks, with each prerequisite completed before the task that needs it.

\smallskip
\noindent{\bf Solution.} Draw an edge from each task to its prerequisites.
Post-order traversal returns each task after its prerequisites. The shared
\code{:base} prerequisite appears only once:

\begtt \hisyntax{clojure}
(def task-prerequisites
  {:app [:lib :config]
   :lib [:base]
   :config [:base]
   :base []})

(mapv :state
      (ss/postorder task-prerequisites [:app]))
;; => [:base :lib :config :app]
\endtt

This ordering is valid because the graph has no dependency cycle. If the
dependency graph has a~cycle, no ordering can put every prerequisite first.

\secc Count tasks after their prerequisites

\noindent{\bf Problem.} A~build plan has four tasks: \code{:app} depends on
\code{:compile} and \code{:docs}; both \code{:compile} and \code{:docs}
depend on \code{:parse}; and \code{:parse} has no prerequisites. For each
task, count the distinct tasks reachable by following its prerequisite
links, including the task itself. The shared \code{:parse} task should only
be counted once for \code{:app}.

\smallskip
\noindent{\bf Solution.} A~post-order \code{:exit} callback runs after a~
task's prerequisites have been processed. Carry the set of prerequisite
tasks already counted, union the sets from completed successors, then count
the resulting set. The callback receives the current visit and its completed
successor visits, and returns fields to add to the current visit:

\begtt \hisyntax{clojure}
(def task-count
  {:exit (fn [{:keys [state]} prerequisites]
           (let [included (into #{}
                                (mapcat #(conj (:included %) (:state %)))
                                prerequisites)]
             {:included (conj included state)
              :size (inc (count included))}))})

(mapv (juxt :state :size)
      (ss/postorder {:app [:compile :docs]
                           :compile [:parse]
                           :docs [:parse]
                           :parse []}
                          [:app]
                          {:extensions [task-count]}))
;; => [[:parse 1] [:compile 2] [:docs 2] [:app 4]]
\endtt

This bottom-up computation handles a~shared prerequisite without counting
it twice in the application's total. Post-order traversal reuses the
completed \code{:parse} result for both dependants; using sets also prevents
counting it twice when their results meet at \code{:app}.

\secc Count ways to break a~word into dictionary words

\noindent{\bf Problem.} Split \code{"catsanddog"} into words from
\code{\{"cat", "cats", "and", "sand", "dog"\}}. How many different
segmentations are possible?

\smallskip
\noindent{\bf Solution.} Use the character offset as a~state. Each
dictionary word that starts at that offset leads to the offset just after
it, so every transition moves forward. A~post-order \code{:exit} callback
can count the segmentations from an offset using the already-counted
successors:

\begtt \hisyntax{clojure}
(def text "catsanddog")
(def dictionary #{"cat" "cats" "and" "sand" "dog"})

(defn next-cuts [start]
  (for [end (range (inc start) (inc (count text)))
        :when (contains? dictionary (subs text start end))]
    end))

(def count-ways
  {:exit (fn [{start :state} suffixes]
           {:ways (if (= start (count text))
                    1
                    (transduce (map :ways) + 0 suffixes))})})

(:ways (last (ss/postorder next-cuts [0]
                                 {:extensions [count-ways]})))
;; => 2
\endtt

Each split point is counted once, while distinct choices of the first word
contribute separate ways to the total.

\secc Count the shortest routes

\noindent{\bf Problem.} Consider a~directed graph with vertices
\code{:a}, \code{:b}, \code{:c}, \code{:d}, \code{:e}, \code{:f},
\code{:g}, and \code{:h}. Its outgoing edges are \code{:a} to \code{:b},
\code{:c}, and \code{:f}; \code{:b} to \code{:d}; \code{:c} to
\code{:d}; \code{:d} to \code{:e}; \code{:f} to \code{:g}; \code{:g} to
\code{:h}; and \code{:h} to \code{:e}. Every edge has cost one. How many
different routes with the fewest edges lead from \code{:a} to \code{:e}?

\smallskip
\noindent{\bf Solution.} First use BFS to find each state's minimum depth.
Keep only edges that advance exactly one depth level; these edges form a~
directed acyclic graph containing all shortest routes. A~post-order
\code{:exit} callback counts the routes from each state to the goal:

\begtt \hisyntax{clojure}
(def route-graph
  {:a [:b :c :f]
   :b [:d]
   :c [:d]
   :d [:e]
   :f [:g]
   :g [:h]
   :h [:e]
   :e []})

(def depth
  (into {} (map (juxt :state :depth))
        (ss/bfs route-graph [:a]
                    {:extensions [ss/depth-extension]})))

(defn shortest-edges [state]
  (filter #(= (inc (depth state)) (depth %))
          (route-graph state)))

(->> (ss/postorder shortest-edges [:a]
                         {:extensions
                          [{:exit (fn [{:keys [state]} successors]
                                    {:count (if (= state :e)
                                              1
                                              (transduce (map :count) + 0
                                                         successors))})}]})
     last
     :count)
;; => 2
\endtt

The two shortest routes share their final edge from \code{:d} to
\code{:e}; the bottom-up count includes both ways of reaching \code{:d}.

\secc Find a~least-cost route with a~negative edge

\noindent{\bf Problem.} A~directed acyclic network has trips
\code{:s -> :a} (4), \code{:s -> :b} (2), \code{:a -> :t} (-3),
\code{:b -> :a} (-1), and \code{:b -> :t} (5). What is the least total
cost from \code{:s} to \code{:t}?

\smallskip
\noindent{\bf Solution.} Dijkstra's priority rule relies on non-negative
edge weights. In a~DAG, post-order traversal can instead evaluate the
cheapest remaining route from each node after its successors are known. The
decoded incoming edge supplies its weight:

\begtt \hisyntax{clojure}
(def dag
  {:s [{:state :a :weight 4} {:state :b :weight 2}]
   :a [{:state :t :weight -3}]
   :b [{:state :a :weight -1} {:state :t :weight 5}]
   :t []})

(def route-cost
  {:exit (fn [{:keys [state]} successors]
           {:cost (if (= state :t)
                    0
                    (apply min
                           (map #(+ (:weight (:via %)) (:cost %))
                                successors)))})})

(:cost (last (ss/postorder dag [:s]
                                 {:decode-edge identity
                                  :extensions [route-cost]})))
;; => -2
\endtt

Topological dynamic programming handles negative weights because the DAG
has no route back to an earlier node. A~general graph with negative cycles
needs a~different algorithm; negative-cycle detection is outside this
library's API.

\secc Find the value of tic-tac-toe under perfect play

\noindent{\bf Problem.} Two players take turns marking an empty square on a~
three-by-three board, with \code{:x} moving first. The first player to
complete a~row, column, or diagonal wins; if the board fills first, the game
is a~draw. Assuming both players choose optimally, what is the game's value
for \code{:x}, where a~win is 1, a~loss is -1, and a~draw is 0?

\smallskip
\noindent{\bf Solution.} Each board is a~state and each legal move adds one
mark, so the game graph is acyclic. An \code{:exit} callback evaluates each
board after its children: \code{:x} chooses the largest child value and
\code{:o} chooses the smallest. Terminal boards receive their value directly:

\begtt \hisyntax{clojure}
(def winning-lines [[0 1 2] [3 4 5] [6 7 8]
                    [0 3 6] [1 4 7] [2 5 8]
                    [0 4 8] [2 4 6]])

(defn board-winner [board]
  (some (fn [[a b c]]
          (when (and (board a) (= (board a) (board b) (board c)))
            (board a)))
        winning-lines))

(defn legal-moves [board]
  (when-not (board-winner board)
    (let [mark (if (= (count (filter #{:x} board))
                      (count (filter #{:o} board)))
                 :x :o)]
      (for [i (range 9) :when (nil? (board i))]
        (assoc board i mark)))))

(def game-value
  {:exit (fn [{board :state} children]
           (let [winner (board-winner board)]
             {:value (cond
                       (= :x winner) 1
                       (= :o winner) -1
                       (empty? children) 0
                       (even? (count (remove nil? board)))
                       (apply max (map :value children))
                       :else (apply min (map :value children)))}))})

(let [positions (ss/postorder legal-moves
                                    [(vec (repeat 9 nil))]
                                    {:extensions [game-value]})]
  [(count positions) (:value (last positions))])
;; => [5478 0]
\endtt

The empty board's value is zero: optimal play from both sides leads to a~
draw. The same bottom-up traversal used for dependency computations can
therefore evaluate a~game tree by minimax.

\secc Schedule ordered preparation steps

\noindent{\bf Problem.} A~lunch needs three preparation steps: chop
vegetables (4~minutes), mix dressing (3~minutes), and plate the food
(2~minutes). They use one work surface, so they must be done in that order.
At what minute does each step begin?

\smallskip
\noindent{\bf Solution.} Treat the lunch as depending on its preparation
steps in work order. A~post-order \code{:next} callback sees the completed
earlier steps in the parent's \code{:previous} field, so it can set the next
step's start time to their total duration:

\begtt \hisyntax{clojure}
(def prep-minutes
  {:chop-vegetables 4
   :mix-dressing 3
   :plate-food 2})

(def prep-schedule
  {:next (fn [parent _transition _visit]
           {:starts-at (reduce + 0
                               (map #(prep-minutes (:state %))
                                    (:previous parent)))})})

(->> (ss/postorder {:lunch [:chop-vegetables
                                  :mix-dressing
                                  :plate-food]
                          :chop-vegetables []
                          :mix-dressing []
                          :plate-food []}
                         [:lunch]
                         {:extensions [prep-schedule]})
     (filter #(contains? % :starts-at)))
;; => ({:state :chop-vegetables :starts-at 0}
;;     {:state :mix-dressing :starts-at 4}
;;     {:state :plate-food :starts-at 7})
\endtt

The callback uses results from earlier siblings to initialize each next
visit. The order matters here because the steps share one work surface.

\secc Restore from the first available copy

\noindent{\bf Problem.} A~document can be restored from a~local copy or,
if that is missing, from a~slower network copy. Check the local copy first.
Which source is used, and can the network lookup be skipped when the local
copy is available?

\smallskip
\noindent{\bf Solution.} Give the preferred source first in the successor
order. Each source reports whether it is available; \code{:more?} lets the
restore visit stop requesting successors once a~completed source succeeds:

\begtt \hisyntax{clojure}
(def available-copies #{:local-copy})

(def restore-extension
  {:more? (fn [{:keys [state]} previous]
            (or (not= state :restore)
                (not-any? :available? previous)))
   :exit (fn [{:keys [state]} successors]
           (if (= state :restore)
             {:source (:state (some #(when (:available? %) %) successors))}
             {:available? (contains? available-copies state)}))})

(ss/postorder {:restore [:local-copy :network-copy]
               :local-copy []
               :network-copy []}
              [:restore]
              {:extensions [restore-extension]})
;; => ({:state :local-copy, :available? true}
;;     {:state :restore, :source :local-copy})
\endtt

The network-copy visit is absent because the local copy succeeded. If it
were unavailable, traversal would continue to the network copy.

\secc Find the edge that closes a~dependency cycle

\noindent{\bf Problem.} A~deployment has two tasks, \code{:api} and
\code{:database}. The API requires the database, and the database requires
the API. Starting at \code{:api}, report each prerequisite edge and indicate
whether it points back to a~task that is still being explored.

\smallskip
\noindent{\bf Solution.} Post-order traversal terminates even when the
graph has cycles. Its \code{:exit} callback receives each successor's
\code{decoded} incoming edge and a~\code{:back-edge} flag, which identifies an
edge to a~task still being explored:

\begtt \hisyntax{clojure}
(def edge-report
  {:exit (fn [_visit successors]
           {:edges (mapv (fn [{:keys [state via back-edge]}]
                           [state (:kind via) back-edge])
                         successors)})})

(ss/postorder {:api [[:database :requires]]
               :database [[:api :requires]]}
              [:api]
              {:decode-edge (fn [[state kind]]
                              {:state state :kind kind})
               :extensions [edge-report]})
;; => ({:state :database :edges [[:api :requires true]]}
;;     {:state :api :edges [[:database :requires false]]})
\endtt

The edge from \code{:database} back to \code{:api} closes the cycle. The
returned order is a~depth-first post-order, not a~valid prerequisite order
when the graph contains a~cycle.

\sec Combining information until it stabilizes

Finally, let information flow between states instead of computing each
state once. First combine service capabilities, then revisit program blocks
when a~loop brings new information.

\secc Find capabilities required by a~service

\noindent{\bf Problem.} A~storefront depends on a~catalog API and an account
API. The catalog API depends on a~database that provides \code{:sql}; the
account API depends on an identity provider that provides \code{:oauth}.
Which capabilities are required by each service, including the storefront?

\smallskip
\noindent{\bf Try it.} Draw edges from each dependency to the service that
uses it. Start with the capabilities provided by the database and identity
provider, then determine how to combine capabilities when dependencies meet.

\smallskip
\noindent{\bf Solution.} Use sets for capabilities and set union when
contributions meet. An adjacency map describes how capabilities flow from
dependencies to services. Seed the database and identity provider with their
capabilities; \code{:bottom \#\{\}} gives newly reached services an empty set,
and \code{:join set/union} combines incoming capabilities. Take the last
progress record's \code{:values} map for the final result:

\begtt \hisyntax{clojure}
(require '[clojure.set :as set])

(def service-dependencies
  {:catalog-db       [:catalog-api]
   :identity-provider [:account-api]
   :catalog-api      [:storefront]
   :account-api      [:storefront]
   :storefront       []})

(def values
  (:values
   (last (ss/propagate service-dependencies
                           {:catalog-db #{:sql}
                            :identity-provider #{:oauth}}
                           {:bottom #{}
                            :join set/union}))))
values
;; => {:catalog-db #{:sql}
;;     :identity-provider #{:oauth}
;;     :catalog-api #{:sql}
;;     :account-api #{:oauth}
;;     :storefront #{:sql :oauth}}
\endtt

The two branches bring different requirements to the storefront, where set
union combines them. Each service's value is the set of capabilities it
inherits from its dependencies. The data is deliberately small; the same
model can be used to inspect larger dependency graphs.

\secc Analyze definitions around a~loop

\noindent{\bf Problem.} A~program has an entry block, a~loop header, a~loop
body, and an exit. The entry assigns to \code{x}; the body assigns to
\code{x} again before returning to the header. At each block, determine
which assignments to \code{x} may reach it. A~definition is a~pair of a~
variable and its assignment label.

\smallskip
\noindent{\bf Try it.} Start with no definitions at the entry. At each
block, remove definitions it overwrites and add definitions it generates.
When paths meet, combine the possible definitions. What changes when the
loop's definition travels back to the header?

\smallskip
\noindent{\bf Solution.} Represent each block by its successors and its
\code{gen}/\code{kill} sets. The transfer function applies those sets before
sending a~contribution along every outgoing edge. Here \code{:project identity}
passes both the block's \code{:state} and its current \code{:value} to the
function. It returns transition maps directly, so \code{:decode-edge identity}
keeps each destination and its contributed value:

\begtt \hisyntax{clojure}
(def blocks
  {:entry  {:successors [:header] :gen #{[:x :initial]} :kill #{}}
   :header {:successors [:body :exit] :gen #{} :kill #{}}
   :body   {:successors [:header]
            :gen #{[:x :loop]}
            :kill #{[:x :initial] [:x :loop]}}
   :exit   {:successors [] :gen #{} :kill #{}}})

(defn transfer [{:keys [state value]}]
  (let [{:keys [successors gen kill]} (blocks state)
        outgoing (into gen (remove kill value))]
    (map (fn [target] {:state target :value outgoing}) successors)))

(def progress
  (ss/propagate transfer
                    {:entry #{}}
                    {:project identity
                     :decode-edge identity
                     :bottom #{}
                     :join set/union}))

(:values (last progress))
;; => {:entry #{}
;;     :header #{[:x :initial] [:x :loop]}
;;     :body #{[:x :initial] [:x :loop]}
;;     :exit #{[:x :initial] [:x :loop]}}
\endtt

The body kills the entry definition and generates the loop definition on
each pass. At the header, contributions from the entry path and the back
edge are joined, so the header's value changes and it is processed again.
The stored sets can only gain definitions from this finite program, so
eventually no block's incoming set changes and processing stops. The
\code{:values} above are the definitions reaching each block before its
transfer runs, not the definitions leaving it. This termination argument
depends on the chosen values and transfer rules; arbitrary propagation
problems need not stabilize.

\bye
